\clearpage
\VerdierPhysicalPageAnchor{0201}%
\begingroup\small\setlength{\parskip}{0pt}\setlength{\emergencystretch}{1.5em}
\setlength{\abovedisplayskip}{4pt plus 2pt minus 2pt}%
\setlength{\belowdisplayskip}{4pt plus 2pt minus 2pt}%

\VerdierSemanticAnchor{ch3d-2-3-3-proof-n0-epimorphism}
where $Z^{-q}$ belongs to $M'$ and the $L^j$, $-q+1\leq j\leq n_0$,
belong to $M$. Consequently, there is an epimorphism
\[
\alpha\colon L^{-q}\longrightarrow Z^{-q},
\qquad\text{where }L^{-q}\in M.
\]
Set
\[
X_{n_0,q}=X_{n_0,q-1}(\alpha,-q),
\qquad
v_{n_0,q}=p(\alpha,-q)\quad\text{(2.1.1)}.
\]

\VerdierSemanticAnchor{ch3d-2-3-3-proof-n0-completion}
It is clear from the construction that properties
$I)_{n_0,q}$ and $III)_{n_0,q}$ hold. Property
$II)_{n_0,q}$ follows from the fact that the set $M'$ has
properties $(E_1)$ and $(E_4)$. We may therefore continue the construction
of the $X_{n_0,i}$. Setting $X_{n_0}=\varprojlim_i X_{n_0,i}$ and taking
the canonical morphism $u_{n_0}\colon X_{n_0}\longrightarrow Y$, we
obtain a complex and a morphism having properties $1)_{n_0}$,
$2)_{n_0}$, $3)_{n_0}$, and $4)_{n_0}$.

\VerdierSemanticAnchor{ch3d-2-3-3-proof-outer-induction}
To construct the complexes $X_i$ and the morphisms
$u_i\colon X_i\longrightarrow X_{i-1}$, we proceed by induction. Suppose
that, for every integer $i$ with $n_0\leq i<q$, we have constructed
complexes $X_i$ and morphisms $u_i\colon X_i\longrightarrow X_{i-1}$
having properties $1)_i$, $2)_i$, $3)_i$, $4)_i$, and $5)_i$,
$n_0\leq i<q$. We must construct a complex $X_q$ and a morphism
$u_q\colon X_q\longrightarrow X_{q-1}$ having properties $1)_q$,
$2)_q$, $3)_q$, $4)_q$, and $5)_q$. First, if $Y^q\in M$, set
$X_q=X_{q-1}$ and $u_q=\operatorname{id}$. Suppose that $Y^q$ does not
belong to $M$, i.e.\ that $M\text{-}\dim(Y^q)>0$.

\VerdierSemanticAnchor{ch3d-2-3-3-proof-q-inner-induction}
To construct the complex $X_q$ and the morphism
$u_q\colon X_q\longrightarrow X_{q-1}$, we proceed once more by
induction. Let $r$ be an integer such that
\[
-q<r\leq -q+M\text{-}\dim(Y^q).
\]
Suppose that we have determined a sequence of complexes $X_{q,l}$,
$-q<l<r$, and morphisms
$v_{q,l}\colon X_{q,l}\longrightarrow X_{q,l-1}$ (we put
$X_{q,-q}=X_{q-1}$) such that:

\VerdierSemanticAnchor{ch3d-2-3-3-proof-q-r-conditions}
\noindent$I')_l$
$(X_{q,l})^j=(X_{q,l-1})^j$ for every $j\neq-l,-l+1$.

\noindent$II')_l$
$(X_{q,l})^{-l}\in M'$ and
$M\text{-}\dim((X_{q,l})^{-l})=M\text{-}\dim(Y^q)-l-q$.

\noindent$III')_l$
$(X_{q,l})^{-l+1}\in M$.

\noindent$IV')_l$
$v_{q,l}$ is an elementary morphism of height $-l+1$ and
$(v_{q,l})^j=\operatorname{id}$ for $j\neq-l,-l+1$.

\VerdierSemanticAnchor{ch3d-2-3-3-proof-q-rminus1-complex}
The complex $X_{q,r-1}$ then has the form
\[
\cdots\longrightarrow L^{-r}\longrightarrow Z^{-r+1}
\longrightarrow L^{-r+2}\longrightarrow\cdots\longrightarrow L^q
\longrightarrow Y^{q+1}\longrightarrow Y^{q+2}\longrightarrow\cdots,
\]
where the $L^i$ belong to $M$ and $Z^{-r+1}$ belongs to $M'$ and satisfies
$M\text{-}\dim(Z^{-r+1})=M\text{-}\dim(Y^q)-r+1-q$. Consequently,
there is an epimorphism
\endgroup
